% Documentación — Portal de Matrícula MPAZ 2027 (maqueta)
% Documento técnico simplificado. Compilar: pdflatex modulo-matricula.tex (2 pasadas).
\documentclass[11pt,a4paper]{article}
\usepackage[spanish,es-tabla]{babel}
\usepackage[utf8]{inputenc}
\usepackage[T1]{fontenc}
\usepackage{lmodern}
\usepackage{microtype}
\usepackage{geometry}
\usepackage{array}
\usepackage{longtable}
\usepackage{booktabs}
\usepackage{enumitem}
\usepackage{fancyhdr}
\usepackage{graphicx}
\usepackage{xcolor}
\usepackage{titlesec}
\usepackage{hyperref}
\geometry{margin=2.5cm}

% Paleta institucional
\definecolor{mpazAzul}{HTML}{0F2537}
\definecolor{mpazDorado}{HTML}{D97706}
\definecolor{mpazFondo}{HTML}{F8FAFC}
\definecolor{mpazBorde}{HTML}{E2E0F0}
\definecolor{mpazGris}{HTML}{64748B}

\hypersetup{colorlinks=true,linkcolor=mpazAzul,urlcolor=mpazAzul,citecolor=mpazAzul,
  pdftitle={Portal de Matrícula MPAZ 2027 — Maqueta},
  pdfauthor={Escuela Escritora Marcela Paz}}

% Encabezado y pie (mismo estilo que el manual de usuarios)
\pagestyle{fancy}
\fancyhf{}
\fancyhead[L]{\small\color{mpazGris}MPAZ 2027 · Portal de Matrícula}
\fancyhead[R]{\small\color{mpazGris}Documentación técnica}
\fancyfoot[C]{\small\color{mpazGris}\thepage}
\renewcommand{\headrulewidth}{0.7pt}
\renewcommand{\headrule}{{\color{mpazDorado}\hrule width\headwidth height\headrulewidth}}

% Títulos
\titleformat{\section}
  {\Large\bfseries\color{mpazAzul}}{\thesection.}{0.6em}{}
  [{\color{mpazDorado}\titlerule[0.8pt]}]
\titleformat{\subsection}
  {\large\bfseries\color{mpazAzul}}{\thesubsection}{0.6em}{}
\titlespacing*{\section}{0pt}{1.6em}{0.8em}

\setlength{\parskip}{0.5em}
\setlength{\parindent}{0pt}
\setlength{\emergencystretch}{2em}
\renewcommand{\arraystretch}{1.2}

% Código en línea: tipografía monoespaciada con color sobrio.
\newcommand{\codigo}[1]{\texttt{\color{mpazAzul}#1}}

% Cajas de nota y de pendiente (robustas, sin paquetes extra).
\newcommand{\caja}[3]{%
  \par\noindent\fcolorbox{#1}{mpazFondo}{%
    \parbox{\dimexpr\linewidth-2\fboxsep-2\fboxrule\relax}{%
      \small\textbf{\color{#1}#2}\ #3}}\par}
\newcommand{\nota}[1]{\caja{mpazAzul}{Nota:}{#1}}
\newcommand{\pendiente}[1]{\caja{mpazDorado}{Pendiente colegio:}{#1}}

\begin{document}
\thispagestyle{empty}
\begin{center}
\vspace*{2cm}
\IfFileExists{../mpaz-matricula/MarcelaPaz.Web/wwwroot/images/logo-mpaz-con-lema.png}{%
  \includegraphics[width=0.18\textwidth]{../mpaz-matricula/MarcelaPaz.Web/wwwroot/images/logo-mpaz-con-lema.png}\par}{}
\vspace{1em}
{\Huge\bfseries\color{mpazAzul} Portal de Matrícula\par}
\vspace{0.4em}
{\large\color{mpazGris} Escuela Escritora Marcela Paz 2027\par}
{\large\color{mpazGris} Documentación del proyecto (maqueta)\par}
\vspace{0.6em}
{\color{mpazDorado}\rule{0.3\textwidth}{1.2pt}}\par
\vspace{0.8em}
Escuela MPAZ\par
2026
\end{center}
\newpage
\tableofcontents
\newpage

\section{Descripción general}
El portal gestiona el ciclo completo de matrícula: búsqueda de estudiante,
ficha en 8 pasos, documentos, firma, validación de Dirección, certificados y
reportes. Incluye salud y medicación TENS, alertas por área y administración
de cuentas. La base de datos es única (\codigo{mpaz}) y contiene el esquema
principal \codigo{mpaz\_matricula} más \codigo{mpaz\_auth} (Identity) y
\codigo{mpaz\_asistencia}.

\section{Arquitectura del código}
Código en \codigo{mpaz-matricula/}, solución
\codigo{MarcelaPaz.Matriculas.slnx} (.NET 10).

\begin{itemize}[leftmargin=*,itemsep=0.25em]
\item \textbf{MarcelaPaz.Web}: controladores, vistas Razor, estilos y servicios.
\item \textbf{MarcelaPaz.Domain}: entidades.
\item \textbf{MarcelaPaz.Infrastructure}: persistencia y documentos PDF.
\end{itemize}

Controladores: \codigo{Ficha}, \codigo{Alumnos}, \codigo{Funcionarios},
\codigo{Autorizacion}, \codigo{Certificados}, \codigo{Reportes},
\codigo{Alertas}, \codigo{Configuracion}, \codigo{Cuenta} y \codigo{Home}.

\section{Base de datos}
Motor PostgreSQL 18 (Neon \codigo{mpaz-database}, región
\codigo{aws-us-east-2}). La instalación nueva se hace en 3 partes, en orden:
\codigo{01-estructura-bd.sql} $\rightarrow$
\codigo{02-programacion-pl.sql} $\rightarrow$
\codigo{03-datos-iniciales.sql}. El mantenimiento usa los pasos \codigo{01} y
\codigo{03}--\codigo{15} (más \codigo{00} rol y \codigo{02} Identity en
\codigo{comun/}). Aplicador: \codigo{aplicar.sh}.

Contenido por pasos: 01 DDL de 27 tablas; 03 matriz de permisos de los 9
roles; 05 año escolar; 08 certificados; 10--11 cursos; 12--13 alertas;
14 roles de Johan; 15 UTP, CONVIVENCIA y TENS.

\section{Roles y permisos}
Nueve roles: SECRETARIA, DIRECCIÓN, INSPECTORÍA, PIE, TENS, UTP, CONVIVENCIA,
DOCENTE y ADMIN\_TI. La matriz vigente está en
\codigo{03-permisos-iniciales.sql}, ajustada por los pasos 14 y 15 (Johan,
17-11-2026).

\begin{itemize}[leftmargin=*,itemsep=0.15em]
\item SECRETARIA crea, edita, formaliza, firma, confirma y anula.
\item DIRECCIÓN valida; INSPECTORÍA corrige datos.
\item PIE, UTP, CONVIVENCIA y DOCENTE solo leen (UTP además exporta).
\item TENS ve salud/medicación y registra dosis.
\item ADMIN\_TI gestiona cuentas y roles.
\end{itemize}

\section{Flujos principales}
\begin{enumerate}[leftmargin=*,itemsep=0.3em]
\item \textbf{Matrícula}: Buscar $\rightarrow$ Nueva/Renovar $\rightarrow$
  pasos 1--8 $\rightarrow$ documentos $\rightarrow$ firma $\rightarrow$ confirmación.
\item \textbf{Validación}: Dirección valida la matrícula y la documentación.
\item \textbf{Certificados}: folio único y hash; la anulación quita validez.
\item \textbf{TENS}: \codigo{RegistrarDosis} con orden VIGENTE y autorización.
\item \textbf{Reportes}: pantalla y CSV según EXPORTAR o LEER por área.
\item \textbf{Cuentas}: ADMIN\_TI crea, desactiva, cierra rol, restablece y desbloquea.
\item \textbf{Web service}: guarda primero en Neon y luego
  \codigo{POST /api/registro/}.
\end{enumerate}

\section{Seguridad y paleta}
Se usan claims \codigo{permiso} (\codigo{AREA:ACCION}) y rol; el menú solo muestra lo
autorizado.

Paleta institucional: azul \codigo{\#0f2537}, dorado
\codigo{\#f59e0b}/\codigo{\#d97706}, fondo \codigo{\#f8fafc} y blanco.

\section{Despliegue y operación}
En desarrollo, el portal corre en \codigo{http://localhost:5301}, con psql local y
user-secrets para \codigo{WebService:Token}. Ramas FT: \codigo{main} intacta y
trabajo en \codigo{mpaz-matricula}.

\end{document}
